\documentclass[11pt,a4paper]{article}
\usepackage[utf8]{inputenc}
\usepackage[english]{babel}
\usepackage{amsmath, amssymb, amsthm}
\usepackage{graphicx}
\usepackage{hyperref}
\usepackage{geometry}
\geometry{margin=2.5cm}

\title{Decoupled Observer Patch Holography (dOPH)}
\author{v04226079-sketch}
\date{April 27, 2026}

\begin{document}

\maketitle

\begin{abstract}
We present a new theoretical model called \textbf{Decoupled Observer Patch Holography (dOPH)}. 
The key parameter $P$ is strictly fixed to the theoretical value 
\[
P \equiv \sqrt{\frac{8}{3}} \approx 1.63299.
\]
Observer decoupling is implemented via a strict projection operator.

The model demonstrates high robustness in extensive synthetic tests on 20,000 galaxies (average stability 98.8\%, LSB galaxies 98.7\%) and successfully reproduces the real Milky Way rotation curve (Eilers et al. 2019), achieving a reduced $\chi^2 = 1.84$ with the fixed $P$. 

A draft fitting pipeline `sparc_fittings_draft_v4.py` has been prepared for the full SPARC dataset (175 galaxies). The project is currently in an active draft stage with all key theoretical components fixed.
\end{abstract}

\section{Current Status as of April 27, 2026}

As of April 27, 2026, the following results have been achieved within the dOPH project:

\begin{itemize}
    \item The parameter $P$ has been fixed to its theoretical value:
    \[
    P \equiv \sqrt{\frac{8}{3}} \approx 1.63299
    \]
    \item Decoupling of the observer is implemented via a strict projection operator.
    \item Extensive synthetic tests were performed on 20,000 galaxies, yielding an average model stability of \textbf{98.8\%} (minimum 98.0\%), and \textbf{98.7\%} for low-surface-brightness (LSB) galaxies.
    \item An independent test was conducted on the real Milky Way rotation curve (Eilers et al. 2019). A reduced $\chi^2 = \textbf{1.84}$ was obtained with the fixed value of $P$.
    \item A draft fitting script `sparc_fittings_draft_v4.py` has been prepared for the full SPARC dataset (175 galaxies). The file \texttt{Rotmod\_LTG.zip} has been placed in the \texttt{data/sparc/} directory.
\end{itemize}

The model demonstrates high robustness on synthetic data and acceptable performance on the Milky Way test. 
The transition to real observational data from SPARC is technically prepared, although the full fitting has not yet been executed due to current computational environment limitations. 
Further work will continue once a more suitable computing environment becomes available.

\section{Conclusion and Outlook}

The dOPH model with the theoretically fixed parameter $P = \sqrt{8/3}$ shows promising results on both synthetic tests and real galactic data. 
Future work includes completing the full SPARC analysis, refining the projection operator, and preparing a preprint for arXiv.

The project remains in an active draft stage, with all core theoretical components (P, projection operator, observer decoupling) firmly established.

\end{document}
